\RequirePackage[l2tabu, orthodox]{nag}
\RequirePackage{silence}
\WarningFilter{nag}{There is no environment ``centering'' }%nag complains because beamer titlepage uses a centering environment
\WarningFilter{nag}{1 complaints in total}
\WarningFilter{ifpdf}{Someone has redefined \pdfoutput}
\WarningFilter{fmtcount}{\ordinal already defined use \FCordinal instead}
\documentclass[french,english]{beamer}
%INSTALL

%avoids a warning
\usepackage[log-declarations=false]{xparse}
\usepackage{fontspec} %font selecting commands
\usepackage{xunicode}
\usepackage{ucharclasses}
%warn about missing characters
\tracinglostchars=2

%REDAC
\usepackage{booktabs}
\usepackage{calc}
\usepackage{tabularx}

\usepackage{etoolbox} %for addtocmd, newtoggle commands
\newtoggle{LCpres}
\toggletrue{LCpres}

\usepackage{mathtools} %load this before babel!
	\mathtoolsset{showonlyrefs,showmanualtags}

\usepackage{babel}
%suppresses the warning about frenchb not modifying the captions (“—” to “:” in “Figure 1 – Legend”).
	\frenchbsetup{AutoSpacePunctuation=false,SuppressWarning=true}

%\usepackage[super]{nth}%better use fmtcount! (loaded by datetime anyway; see below about pbl with warnings and package silence)
\usepackage{listings} %typeset source code listings
	\lstset{language=XML,tabsize=2,captionpos=b,basicstyle=\NoAutoSpacing}%NoAutoSpacing avoids space before colon or ?}%,literate={"}{{\tt"}}1, keywordstyle=\fontspec{Latin Modern Mono Light}\textbf, emph={String, PreparedStatement}, emphstyle=\fontspec{Latin Modern Mono Light}\textbf, language=Java, basicstyle=\small\NoAutoSpacing\ttfamily, aboveskip=0pt, belowskip=0pt, showstringspaces=false
\usepackage[nolist,smaller,printonlyused]{acronym}%,smaller option produces warnings from relsize in some cases, it seems.% Note silence and acronym and hyperref make (xe)latex crash when ac used in section (http://tex.stackexchange.com/questions/103483/strange-packages-interaction-acronyms-silence-hyperref), rather use \section{\texorpdfstring{\acs{UE}}{UE}}.
\usepackage{fmtcount}
\usepackage[nodayofweek]{datetime}%must be loaded after the babel package. However, loading it after {nth} generates a warning from fmtcount about ordinal being already defined. Better load it before nth? (then we can remove the silence package which creates possible crashes, see above.) Or remove nth?
%\usepackage{xspace}%do we need this?
\nottoggle{LCpres}{
	\usepackage[textsize=small]{todonotes}
}
\iftoggle{LCpres}{
	%remove pdfusetitle (implied by beamer)
	\usepackage{hyperref}
}{
% option pdfusetitle must be introduced here, not in hypersetup.
	\usepackage[pdfusetitle]{hyperref}
}
%breaklinks makes links on multiple lines into different PDF links to the same target.
%colorlinks (false): Colors the text of links and anchors. The colors chosen depend on the the type of link. In spite of colored boxes, the colored text remains when printing.
%linkcolor=black: this leaves other links in colors, e.g. refs in green, don't print well.
%pdfborder (0 0 1, set to 0 0 0 if colorlinks): width of PDF link border
%hidelinks
\hypersetup{breaklinks, bookmarksopen}
%add hyperfigures=true in hypersetup (already defined in article mode)
\iftoggle{LCpres}{
	\hypersetup{hyperfigures}
}

%in Beamer, sets url colored links but does not change the rest of the colors (http://tex.stackexchange.com/questions/13423/how-to-change-the-color-of-href-links-for-real)
%\hypersetup{breaklinks,bookmarksopen,colorlinks=true,urlcolor=blue,linkcolor=,hyperfigures=true}
% hyperref doc says: Package bookmark replaces hyperref’s bookmark organization by a new algorithm (...) Therefore I recommend using this package.
\usepackage{bookmark}

% center floats by default, but do not use with float
% \usepackage{floatrow}
% \makeatletter
% \g@addto@macro\@floatboxreset\centering
% \makeatother
\nottoggle{LCpres}{
	\usepackage{enumitem} %follow enumerate by a string saying how to display enumeration
}
\usepackage{ragged2e} %new com­mands \Cen­ter­ing, \RaggedLeft, and \RaggedRight and new en­vi­ron­ments Cen­ter, FlushLeft, and FlushRight, which set ragged text and are eas­ily con­fig­urable to al­low hy­phen­ation (the cor­re­spond­ing com­mands in LaTeX, all of whose names are lower-case, pre­vent hy­phen­ation al­to­gether). 
\usepackage{siunitx} %[expproduct=tighttimes, decimalsymbol=comma] ou (plus récent ?) [round-mode=figures, round-precision=2, scientific-notation = engineering]
\sisetup{detect-all, locale = FR, strict}% to detect e.g. when in math mode (use a math font) - check whether this makes sense with strict
\usepackage{braket} %for \Set
\usepackage{natbib}
\usepackage{doi}

\usepackage{amsmath,amsthm}
% \usepackage{amsfonts} %not required?
% \usepackage{dsfont} %for what?
%unicode-math overwrites the following commands from the mathtools package: \dblcolon, \coloneqq, \Coloneqq, \eqqcolon. Using the other colon-like commands from mathtools will lead to inconsistencies. Plus, Using \overbracket and \underbracke from mathtools package. Use \Uoverbracket and \Uunderbracke for original unicode-math definition.
%use exclusively \mathbf and choose math bold style below.
\usepackage[warnings-off={mathtools-colon, mathtools-overbracket}, bold-style=ISO]{unicode-math}

\defaultfontfeatures{
	Fractions=On,
	Mapping=tex-text% to turn "--" into dashes, useful for bibtex%%
}
%\defaultfontfeatures[\rmfamily]{
%	Fractions=On,
%	Mapping=% to leave " alone (disable the default mapping tex-text; requires loading the font afterwards?)
%}
\newfontfamily\xitsfamily{XITS}
\newfontfamily\texgyretermesfamily{TeX Gyre Termes}
\newfontfamily\texgyreherosfamily{TeX Gyre Heros}
\newfontfamily\lmfamily{Latin Modern Roman}
\iftoggle{LCpres}{
	\setmainfont{Latin Modern Roman}
	\setsansfont{Latin Modern Sans}
	\setmonofont{Latin Modern Mono}
}{
	\setmainfont{TeX Gyre Termes}
	\setsansfont{TeX Gyre Heros}
	\setmonofont{TeX Gyre Cursor}
}
\defaultfontfeatures{}%disable default font features to avoid warnings with math fonts.
\setmathfont{XITS Math}
\setmathfont[range={\mathcal,\mathbfcal},StylisticSet=1]{XITS Math}
%silently delays restoring normal font when using e.g. “≠)” or “≠ other word”.
%\setTransitionTo{Arrows}{\xitsfamily}
%\setTransitionTo{MathematicalOperators}{\xitsfamily}
%this solution explicitly fails on “≠)” and works for “≠ other word.”
\setTransitionsFor{MathematicalOperators}{\begingroup\xitsfamily}{\endgroup}
\setTransitionsFor{Arrows}{\begingroup\xitsfamily}{\endgroup}
%texgyretermes does not have ✓, XITS does.
\setTransitionsFor{Dingbats}{\begingroup\xitsfamily}{\endgroup}
%this also works, but it is more complex
%\def\ResetTransitionTo#1{%
%  \XeTeXinterchartoks 255 \csname#1Class\endcsname{\relax}}
%\setTransitionsFor{MathematicalOperators}
%  {\begingroup\ResetTransitionTo{MathematicalOperators}\xitsfamily}
%  {\endgroup}

\usepackage{cleveref}% cleveref should go "laster" than hyperref
%GRAPHICS
\usepackage{pgf}
\usepackage{pgfplots}
	\usetikzlibrary{babel, matrix, fit, plotmarks, calc, trees, shapes.geometric, positioning, plothandlers, arrows, shapes.multipart}
\pgfplotsset{compat=1.11}
\usepackage{graphicx}

\graphicspath{{graphics/},{graphics-dm/}}
\DeclareGraphicsExtensions{.pdf}
\LetLtxMacro\SavedIncludeGraphics\includegraphics
\AtBeginDocument{
	\def\includegraphics#1#{% #1 catches optional stuff (star/opt. arg.)
		\IncludeGraphicsAux{#1}%
	}%
}
\newcommand*{\IncludeGraphicsAux}[2]{%
	\XeTeXLinkBox{%
		\SavedIncludeGraphics#1{#2}%
	}%
}%

%HACKING
\usepackage{printlen}
\uselengthunit{mm}
% 	\newlength{\templ}% or LenTemp?
% 	\setlength{\templ}{6 pt}
% 	\printlength{\templ}
\usepackage{scrhack}% load at end. Corrects a bug in float package, which is outdated but might be used by other packages
\usepackage{xltxtra} %somebody said that this is loaded by fontspec, but does not seem correct: if not loaded explicitly, does not appear in the log and \showhyphens is not corrected.

%Beamer-specific
%do not remove babel, which beamer uses (beamer uses the \translate command for the appendix); but french can be removed.
\iftoggle{LCpres}{
	\usepackage{appendixnumberbeamer}
	\setbeamertemplate{navigation symbols}{} 
	\usepackage{preamble/beamerthemeParisFrance}
	\usefonttheme{professionalfonts}
	\setcounter{tocdepth}{10}
	%From: http://tex.stackexchange.com/questions/168057/beamer-with-xelatex-on-texlive2013-enumerate-numbers-in-black
%I don’t think it’s useful to submit this as a bug: nothing has been solved since March, 2015. See: https://bitbucket.org/rivanvx/beamer/issues?status=resolved.

\setbeamertemplate{enumerate item}
{
  \begin{pgfpicture}{-1ex}{-0.65ex}{1ex}{1ex}
    \usebeamercolor[fg]{item projected}
    {\pgftransformscale{1.75}\pgftext{\normalsize\pgfuseshading{bigsphere}}}
    {\pgftransformshift{\pgfpoint{0pt}{0.5pt}}
      \pgftext{\usebeamercolor[fg]{item projected}\usebeamerfont*{item projected}\insertenumlabel}}
  \end{pgfpicture}%
}

\setbeamertemplate{enumerate subitem}
{
  \begin{pgfpicture}{-1ex}{-0.55ex}{1ex}{1ex}
    \usebeamercolor[fg]{subitem projected}
    {\pgftransformscale{1.4}\pgftext{\normalsize\pgfuseshading{bigsphere}}}
    \pgftext{%
      \usebeamercolor[fg]{subitem projected}%
      \usebeamerfont*{subitem projected}%
      \insertsubenumlabel}
  \end{pgfpicture}%
}

\setbeamertemplate{enumerate subsubitem}
{
  \begin{pgfpicture}{-1ex}{-0.55ex}{1ex}{1ex}
    \usebeamercolor[fg]{subsubitem projected}
    {\pgftransformscale{1.4}\pgftext{\normalsize\pgfuseshading{bigsphere}}}
    \pgftext{%
      \usebeamercolor[fg]{subsubitem projected}%
      \usebeamerfont*{subitem projected}%
      \insertsubsubenumlabel}
  \end{pgfpicture}%
}


}
% \newcommand{\citep}{\cite}%Better: leave natbib.
% \setbeamersize{text margin left=0.1cm, text margin right=0.1cm} 
% \usetheme{BrusselsBelgium}%no, replace with paris
%\usetheme{ParisFrance}, no, usepackage better!
% Tex Gyre takes too much space, replace with Latin Modern Roman / Sans / Mono.
% Difference when loading explicitly Latin Modern Sans (compared to not using \setsansfont at all):
% the font LMSans17-Regular appears in the document;
% the title of the slides appears differently;
% it does not say (in the log file):
% > LaTeX Font Info:    Font shape `EU1/lmss/m/it' in size <10.95> not available
% > (Font)              Font shape `EU1/lmss/m/sl' tried instead on input line 85.
% > LaTeX Font Info:    Try loading font information for EU1+lmtt on input line 85.

%tikzposter-specific
%remove \usepackage{ragged2e}: causes 1=1 to be printed in the middle of the poster. (Anyway prints a warning about those characters being missing.)
%put [french, english] options next to \usepackage{babel} to avoid warning of 

\newcommand{\R}{ℝ}
\newcommand{\N}{ℕ}
\newcommand{\Z}{ℤ}
\newcommand{\card}[1]{\lvert{#1}\rvert}
\newcommand{\powerset}[1]{\mathscr{P}(#1)}%\mathscr rather than \mathcal: scr is rounder than cal (at least in XITS Math).
\newcommand{\suchthat}{\;\ifnum\currentgrouptype=16 \middle\fi|\;}
%\newcommand{\Rplus}{\reels^+\xspace}

\AtBeginDocument{%
	\renewcommand{\epsilon}{\varepsilon}
% we want straight form of \phi for mathematics, as recommended in UTR #25: Unicode support for mathematics.
%	\renewcommand{\phi}{\varphi}
}

% with amssymb, but I don’t want to use amssymb just for that.
% \newcommand{\restr}[2]{{#1}_{\restriction #2}}
%\newcommand{\restr}[2]{{#1\upharpoonright}_{#2}}
\newcommand{\restr}[2]{{#1|}_{#2}}%sometimes typed out incorrectly within \set.
%\newcommand{\restr}[2]{{#1}_{\vert #2}}%\vert errors when used within \Set and is typed out incorrectly within \set.
\DeclareMathOperator*{\argmax}{arg\,max}
\DeclareMathOperator*{\argmin}{arg\,min}


%Voting
\newcommand{\allalts}{\mathscr{A}}
\newcommand{\alts}{A}
\newcommand{\feasalts}{F}
\newcommand{\allvoters}{\mathscr{N}}
\newcommand{\voters}{N}
\newcommand{\allsystems}{\mathcal{G}}
\newcommand{\prof}{\mathbf{R}}
\newcommand{\allprofs}{\mathbfcal{R}}
\newcommand{\linors}{\mathcal{L}(\alts)}

\newcommand{\pbasic}[1]{\prof^{#1}_\epsilon}
\newcommand{\pelem}[1]{\prof^{#1}_e}
\newcommand{\pcycl}[1]{\prof^{#1}_c}
\newcommand{\pcycllong}[1]{\prof^{#1}_{cl}}
\newcommand{\pinv}[1]{\overline{\prof_{#1}}}
\newcommand{\dmap}{{\xitsfamily δ}}
%powerset without zero
\newcommand{\powersetz}[1]{\mathcal{P}_\emptyset(#1)}

%logic atom
%⟼ (long)
\DeclareDocumentCommand{\lato}{ O{\prof} O{\alts} }{[#1 \!⟼\! #2]}
%logic atom in
%↝, \stackrel{\in}{\mapsto}, ➲, ⥹
\newcommand{\tightoverset}[2]{%
  \mathop{#2}\limits^{\vbox to -.5ex{\kern-0.9ex\hbox{$#1$}\vss}}}
\DeclareDocumentCommand{\latoin}{ O{\prof} O{\alpha} }{[#1 \tightoverset{\in}{⟼} #2]}
\newcommand{\alllang}{\mathcal{L}}
\newcommand{\ltru}{\texttt{T}}
\newcommand{\lfal}{\texttt{F}}
\newcommand{\laxiom}[1]{{\texgyreherosfamily{\textsc{#1}}}}

%ARG TH
\newcommand{\AF}{\mathcal{AF}}
\newcommand{\labelling}{\mathcal{L}}
\newcommand{\labin}{\textbf{in}\xspace}
\newcommand{\labout}{\textbf{out}}
\newcommand{\labund}{\textbf{undec}\xspace}
\newcommand{\nonemptyor}[2]{\ifthenelse{\equal{#1}{}}{#2}{#1}}
\newcommand{\gextlab}[2][]{
	\labelling{\mathcal{GE}}_{(#2, \nonemptyor{#1}{\ibeatsr{#2}})}
}
\newcommand{\allargs}{A^*}
\newcommand{\args}{A}
\newcommand{\ar}{a}
\newcommand{\ext}{\mathcal{E}}

%MCDA+Arg
\newcommand{\dm}{d}
\newcommand{\ileadsto}{\rightcurvedarrow}
\newcommand{\mleadsto}[1][\eta]{\rightcurvedarrow_{#1}}
\newcommand{\ibeats}{\vartriangleright}
\newcommand{\mbeats}[1][\eta]{\vartriangleright_{#1}}


%MISC
\newcommand{\lequiv}{\Vvdash}
\newcommand{\weightst}{W^{\,t}}

%MCDA classical
\newcommand{\crits}{\mathcal{J}}

%Sorting
\newcommand{\cats}{\mathcal{C}}
\newcommand{\catssubsets}{2^\cats}
\newcommand{\catgg}{\vartriangleright}
\newcommand{\catll}{\vartriangleleft}
\newcommand{\catleq}{\trianglelefteq}
\newcommand{\catgeq}{\trianglerighteq}
\newcommand{\alttoc}[2][x]{(#1 \xrightarrow{} #2)}
\newcommand{\alttocat}[3]{(#2 \xrightarrow{#1} #3)}
\newcommand{\alttoI}{(x \xrightarrow{} \left[\underline{C_x}, \overline{C_x}\right])}
\newcommand{\alttocatdm}[3][t]{\left(#2 \thinspace \raisebox{-3pt}{$\xrightarrow{#1}$}\thinspace #3\right)}
\newcommand{\alttocatatleast}[2]{\left(#1 \thinspace \raisebox{-3pt}{$\xrightarrow[]{≥}$}\thinspace #2\right)}
\newcommand{\alttocatatmost}[2]{\left(#1 \thinspace \raisebox{-3pt}{$\xrightarrow[]{≤}$}\thinspace #2\right)}

\definecolor{darkgreen}{rgb}{0,0.6,0}
\newcommand{\commentOC}[1]{{\selectlanguage{french}{\todo{OC : #1}}}}
%Or: \todo[color=green!40]
\newcommand{\innote}[1]{{\scriptsize{#1}}}

%this probably requires outdated float package, see doc KomaScript for an alternative.
% \newfloat{program}{t}{lop}
% \floatname{program}{PM}

%definition, theorem, lemma, example environments, qed trickery are only needed in article mode (not Beamer)
\nottoggle{LCpres}{
%style is plain by default (italic text)
	\newtheorem{definition}{Definition}
	\newtheorem{theorem}{Theorem}
%no italic: expected.
%http://tex.stackexchange.com/questions/144653/italicizing-of-amsthm-package
	\newtheorem{lemma}{Lemma}
%\crefname{axiom}{axiom}{axioms}%might be needed for workaround bug in cref when defining new theorems?

%\ifdefined\theorem\else
%\newtheorem{theorem}{\iflanguage{english}{Theorem}{Théorème}}
%\fi

\theoremstyle{remark}
	\newtheorem{examplex}{Example}
	\newtheorem{remarkx}{Remark}

%trickery allowing use of \qedhere and the like.
\newenvironment{example}{
	\pushQED{\qed}\renewcommand{\qedsymbol}{$\triangle$}\examplex
}{
	\popQED\endexamplex
}
\newenvironment{remark}{
	\pushQED{\qed}\renewcommand{\qedsymbol}{$\triangle$}\remarkx
}{
	\popQED\endremarkx
}
}

%which line breaks are chosen: accept worse lines, therefore reducing risk of overfull lines. Default = 200
\tolerance=2000
%accept overfull hbox up to...
\hfuzz=2cm
%reduces verbosity about the bad line breaks
\hbadness 5000
%sloppy sets tolerance to 9999
\apptocmd{\sloppy}{\hbadness 10000\relax}{}{}

\bibliographystyle{abbrvnat}
%or \bibliographystyle{apalike} for presentations?

%doi package uses old-style dx.doi url, see 3.8 DOI system Proxy Server technical details, “Users may resolve DOI names that are structured to use the DOI system Proxy Server (http://doi.org (preferred) or http://dx.doi.org).”, https://www.doi.org/doi_handbook/3_Resolution.html
\makeatletter
\patchcmd{\@doi}{dx.doi.org}{doi.org}{}{}
\makeatother

% WRITING
%\newcommand{\ie}{i.e.\@\xspace}%to try
%\newcommand{\eg}{e.g.\@\xspace}
%\newcommand{\etal}{et al.\@\xspace}
\newcommand{\ie}{i.e.\ }
\newcommand{\eg}{e.g.\ }
\newcommand{\mkkOK}{\checkmark}%\color{green}{\checkmark}
\newcommand{\mkkREQ}{\ding{53}}%requires pifont?%\color{green}{\checkmark}
\newcommand{\mkkNO}{}%\text{\color{red}{\textsf{X}}}

\makeatletter
\newcommand{\boldor}[2]{%
	\ifnum\strcmp{\f@series}{bx}=\z@
		#1%
	\else
		#2%
	\fi
}
\newcommand{\textstyleElProm}[1]{\boldor{\MakeUppercase{#1}}{\textsc{#1}}}
\makeatother
\newcommand{\electre}{\textstyleElProm{Électre}\xspace}
\newcommand{\electreIv}{\textstyleElProm{Électre Iv}\xspace}
\newcommand{\electreIV}{\textstyleElProm{Électre IV}\xspace}
\newcommand{\electreIII}{\textstyleElProm{Électre III}\xspace}
\newcommand{\electreTRI}{\textstyleElProm{Électre Tri}\xspace}
% \newcommand{\utadis}{\texorpdfstring{\textstyleElProm{utadis}\xspace}{UTADIS}}
% \newcommand{\utadisI}{\texorpdfstring{\textstyleElProm{utadis i}\xspace}{UTADIS I}}

%TODO
% \newcommand{\textstyleElProm}[1]{{\rmfamily\textsc{#1}}} 


%const
\newcommand{\tikzboxit}{\path node[draw, overlay, inner sep=0.6mm, fit=(boxed), rectangle] {};}%

\newlength{\GraphsNodeSep}
\setlength{\GraphsNodeSep}{7mm}

% MCDA Drawing Sorting
\newlength{\MCDSCatHeight}
\setlength{\MCDSCatHeight}{6mm}
\newlength{\MCDSAltHeight}
\setlength{\MCDSAltHeight}{4mm}
%separation between two vertical alts
\newlength{\MCDSAltSep}
\setlength{\MCDSAltSep}{2mm}
\newlength{\MCDSCatWidth}
\setlength{\MCDSCatWidth}{3cm}
\newlength{\MCDSEvalRowHeight}
\setlength{\MCDSEvalRowHeight}{6mm}
\newlength{\MCDSAltsToCatsSep}
\setlength{\MCDSAltsToCatsSep}{1.5cm}
\newcounter{MCDSNbAlts}
\newcounter{MCDSNbCats}
\newlength{\MCDSArrowDownOffset}
\setlength{\MCDSArrowDownOffset}{0mm}

\tikzset{/Graphs/dot/.style={
	shape=circle, fill=black, inner sep=0, minimum size=1mm
}}
\tikzset{/MC/D/S/alt/.style={
	shape=rectangle, draw=black, inner sep=0, minimum height=\MCDSAltHeight, minimum width=2.5cm, anchor=north east
}}
\tikzset{MC/D/S/pref/.style={
	shape=ellipse, draw=gray, thick
}}
\tikzset{/MC/D/S/cat/.style={
	shape=rectangle, draw=black, inner sep=0, minimum height=\MCDSCatHeight, minimum width=\MCDSCatWidth, anchor=north west
}}
\tikzset{/MC/D/S/evals matrix/.style={
	matrix, row sep=-\pgflinewidth, column sep=-\pgflinewidth, nodes={shape=rectangle, draw=black, inner sep=0mm, text depth=0.5ex, text height=1em, minimum height=\MCDSEvalRowHeight, minimum width=12mm}, nodes in empty cells, matrix of nodes, inner sep=0mm, outer sep=0mm, row 1/.style={nodes={draw=none, minimum height=0em, text height=, inner ysep=1mm}}
}}
\tikzset{profile matrix/.style={
	matrix of math nodes, column sep=3mm, row sep=2mm, nodes={inner sep=0.5mm, anchor=base}
}}
\tikzset{rank-profile matrix/.style={
	matrix of math nodes, column sep=3mm, row sep=2mm, nodes={anchor=base}, column 1/.style={nodes={inner sep=0.5mm}}, row 1/.style={nodes={inner sep=0.5mm}}
}}
\tikzset{rank-vector/.style={
	draw, rectangle, inner sep=0, outer sep=1mm
}}
\tikzset{isolated rank-vector/.style={
	draw, matrix of math nodes, column sep=3mm, inner sep=0, matrix anchor=base, nodes={anchor=base, inner sep=.33em}, ampersand replacement=\&
}}

% GUI
\tikzset{/GUI/button/.style={
	rectangle, very thick, rounded corners, draw=black, fill=black!40%, top color=black!70, bottom color=white
}}

% Beliefs
\tikzset{/Beliefs/D/S/attacker/.style={
	shape=rectangle, draw, minimum size=8mm
}}
\tikzset{/Beliefs/D/S/supporter/.style={
	shape=circle, draw
}}

\newcommand{\tikzmark}[1]{%
	\tikz[overlay, remember picture, baseline=(#1.base)] \node (#1) {};%
}


\begin{acronym}
\acro{AMCD}{Aide Multicritère à la Décision}
\acro{ASA}{Argument Strength Assessment}
\acro{DA}{Decision Analysis}
\acro{DM}{Decision Maker}
\acro{DPr}{Deliberated Preferences}
\acro{DRSA}{Dominance-based Rough Set Approach}
\acro{DSS}{Decision Support Systems}
\acrodefplural{DSS}{Decision Support Systems}
% \newacroplural{DSS}[DSSes]{Decision Support Systems}
\acro{EJOR}{European Journal of Operational Research}
\acro{LNCS}{Lecture Notes in Computer Science}
\acro{MCDA}{Multicriteria Decision Aid}
\acro{MIP}{Mixed Integer Program}
\acro{NCSM}{Non Compensatory Sorting Model}
\acro{PL}{Programme Linéaire}
\acro{PLNE}{Programme Linéaire en Nombres Entiers}
\acro{PM}{Programme Mathématique}
\acro{MP}{Mathematical Program}
\acro{MIP}{Mixed Integer Program}
% \newacroplural{PM}{Programmes Mathématiques}
%acrodefplural since version 1.35, my debian has \ProvidesPackage{acronym}[2009/01/25, v1.34, Support for acronyms (Tobias Oetiker)]
\acrodefplural{PM}{Programmes Mathématiques}
\acro{PMML}{Predictive Model Markup Language}
\acro{RESS}{Reliability Engineering \& System Safety}
\acro{SMAA}{Stochastic Multicriteria Acceptability Analysis}
\acro{URPDM}{Uncertainty and Robustness in Planning and Decision Making}
\acro{XML}{Extensible Markup Language}
\end{acronym}


\title{Arguing about voting rules}
\subject{Social Choice}
\keywords{axioms, Borda, Condorcet, explanation, argumentation}
\author{Olivier Cailloux}
\institute[LAMSADE]{LAMSADE, Université Paris-Dauphine \\ Joint work with Ulle Endriss, University of Amsterdam}
\date{\formatdate{16}{2}{2017}}
\begin{document}
\begin{frame}[plain]
	\tikz[remember picture,overlay]{
		\path (current page.south west) node[anchor=south west, inner sep=0] {
			\includegraphics[height=1cm]{LAMSADE95.jpg}
		};
		\path (current page.south) ++ (0, 1mm) node[anchor=south, inner sep=0] {
			\includegraphics[height=9mm]{Dauphine.jpg}
		};
		\path (current page.south east) node[anchor=south east, inner sep=0] {
			\includegraphics[height=1cm]{PSL.png}
		};
	}
	\titlepage
\end{frame}
\addtocounter{framenumber}{-1}

\begin{frame}
	\frametitle{Outline}
	\tableofcontents[hideallsubsections, sectionstyle=shaded/show]
\end{frame}

\AtBeginSection{
	\begin{frame}
		\frametitle{Outline}
		\tableofcontents[currentsection, hideallsubsections]
	\end{frame}
}

\section{Introduction to Social Choice}
\subsection{Definition}
\begin{frame}
	\frametitle{\subsecname}
	\begin{definition}[Social Choice]
		Designing and analyzing methods for collective decision making
	\end{definition}
	\begin{itemize}
		\item Focus on: mathematical analysis {\tiny preferably non empirical}
		\item (Usually) no focus on: collective social behavior
	\end{itemize}
\end{frame}

\subsection{The Condorcet paradox}
\begin{frame}
	\frametitle{\subsecname}
	\begin{itemize}
		\item Why should we analyze decision schemes mathematically?
		\item Here is a preliminary argument
		\item 3 alternatives : $a, b, c$
		\item 3 voters : $1, 2, 3$
		\item Compare alternatives pairwise
		\item $a$ VS $b$: 2 votes for $a$
		\item $b$ VS $c$: 2 votes for $b$
		\item Conclusion? \pause $a \succ c$? Not so simple!
	\end{itemize}
	\begin{equation}
		\begin{array}{rrr}
			a	&b	&c\\
			b	&c	&a\\
			c	&a	&b\\
		\end{array}
	\end{equation}
\end{frame}

\subsection{Two key moments}
\begin{frame}
	\frametitle{\subsecname: 1. Borda, Condorcet}
	\begin{itemize}
		\item Nicolas de Condorcet [1743 – 1794]
		\item Jean-Charles de Borda [1733 – 1799]
		\item Start of \emph{mathematical, systematic} analysis of voting schemes
		\item Rich period of constitution writing
		\item Condorcet initially involved in the writing of the French constitution
		\item Borda’s rank-order method was first proposed orally at the French Academy of Science in 1770 \citep{suzumura_introduction_2002}
	\end{itemize}
\end{frame}

\begin{frame}
	\frametitle{\subsecname: 2. Arrow}
	\begin{itemize}
		\item up to Arrow: mostly analysis of individual voting schemes
		\item Kenneth Arrow [1921 – ]
		\item Analysis of \emph{classes} of voting schemes
		\item According to properties
		\item[⇒] {\tiny (for example)} Nobel price in 1998 to Amartya Sen (welfare economics)
	\end{itemize}
\end{frame}

\subsection{This talk VS reality}
\begin{frame}
	\frametitle{\subsecname}
	\begin{itemize}
		\item Elections happen IRL (? \pause in real life) \pause
		\item Quite far from the setting examined here
		\item Preferences expressed in various ways
		\item Procedures may not fit this framework
		\item Social Choice in general relevant to this part of reality
		\item In this talk: \emph{limited} connexion
	\end{itemize}
\end{frame}

\section{Voting rules}
\subsection{Definition}
\begin{frame}
	\frametitle{\subsecname}
	\begin{itemize}
		\item Voting rule: a systematic way of aggregating different opinions and decide
		\item Multiple reasonable ways of doing this
		\item Different voting rules have different interesting properties
		\item None satisfy all desirable properties
	\end{itemize}
\end{frame}

\begin{frame}[fragile]
	\frametitle{Voting rule}
	
	\begin{description}[Alternatives]
		\item[Alternatives] $\allalts = \set{a, b, c, d, \ldots}$
		\item[Voters] $\allvoters = \set{1, 2, \ldots}$
		\item[Profile] function $\prof$ from $\allvoters$ to linear orders on $\allalts$.
		\item[Voting rule] function $f$ mapping each $\prof$ to winners $\emptyset \subset \alts \subseteq \allalts$.
	\end{description}
	\vfill
	\begin{center}
		\begin{tikzpicture}
			\path node[profile matrix] (profile) {
				\prof(1)&
				\prof(2)
				\\
				| (profile11) | a&
				b
				\\
				b&
				a
				\\
				c&
				| (profile32) | c
				\\
			};
			\path ($(profile.south west)!.5!(profile.south east)$) ++ (0, -5mm) node {$\prof$};
			
			\path node[draw, rectangle, fit=(profile11) (profile32), outer xsep=2mm, outer ysep=1mm] (justprofile) {};
			\path (justprofile.east) ++ (2.5cm, 0) node[inner sep=0] (winners) {\mbox{} $\alts = \Set{a, b}$};
			\path[draw, ->] (justprofile.east) to[bend left=35] node[anchor=south] {$f$} (winners.west);

			\path[draw, decorate, decoration={brace, mirror}] (justprofile.south west) -- (justprofile.south east);
		\end{tikzpicture}
	\end{center}
\end{frame}

\subsection{Example of a profile}
\begin{frame}[fragile]
	\frametitle{\subsecname}
	\begin{equation}
		\begin{array}{lrrrrrr}
			&\multicolumn{6}{c}{\text{nb voters}}\\
		\cmidrule{2-7}
				&33	&16	&3	&8	&18	&22	\\
		\midrule
			1	&a	&b	&c	&c	&d	&e	\\
			2	&b	&d	&d	&e	&e	&c	\\
			3	&c	&c	&b	&b	&c	&b	\\
			4	&d	&e	&a	&d	&b	&d	\\
			5	&e	&a	&e	&a	&a	&a	\\
		\end{array}
	\end{equation}
	Who wins?\pause
	\begin{itemize}
		\item Most top-1: $a$
		\item $c$ is in the top 3 for everybody
		\item delete worst first, lowest nb of pref: $c$, $b$, $e$, $a$ ⇒ $d$
		\item delete worst first, from bottom: $a$, $e$, $d$, $b$ ⇒ $c$
		\item Borda: $b$
%		\item Condorcet: $c$
	\end{itemize}
\end{frame}

\subsection{Two voting rules}
\begin{frame}
	\frametitle{Borda}
	
Given a profile $\prof$:
	\begin{itemize}
		\item score of $a \in \allalts$: number of alternatives it beats
		\item the highest scores win
	\end{itemize}
	
	\begin{equation}
		\prof =
		\begin{array}{rrrrr}
			a	&	a	&	a	&	b	&	b\\
			b	&	b	&	b	&	c	&	c\\
			c	&	c	&	c	&	a	&	a
		\end{array}
	\end{equation}
	\begin{itemize}
		\item score $a$ is\dots? \pause $2 + 2 + 2 = 6$
		\item score $b$ is $1 + 1 + 1 + 2 + 2 = 7$
		\item score $c$ is $1 + 1 = 2$
	\end{itemize}
	Winner: $b$.
\end{frame}

\begin{frame}
	\frametitle{Condorcet’s principle}
	\begin{block}{Condorcet’s principle}
		We ought to take the Condorcet winner as sole winner if it exists.
		\begin{itemize}
			\item $a$ \emph{beats} $b$ iff more than half the voters prefer $a$ to $b$.
			\item $a$ is a \emph{Condorcet winner} iff $a$ beats every other alternatives.
		\end{itemize}
	\end{block}
	\vfill
	\begin{equation}
		\prof =
		\begin{array}{rrrrr}
			a	&	a	&	a	&	b	&	b\\
			b	&	b	&	b	&	c	&	c\\
			c	&	c	&	c	&	a	&	a
		\end{array}
	\end{equation}
	 Who wins? \pause $a$
\end{frame}

\begin{frame}
	\frametitle{Condorcet’s principle and a voting rule}
	\begin{itemize}
		\item Condorcet’s principle does not define a voting rule. Why? \pause
		\item No winner is defined when no Condorcet winner
	\end{itemize}
	\begin{equation}
		\prof =
		\begin{array}{rrr}
			b	&	c	&	d\\
			c	&	d	&	b\\
			a	&	b	&	a\\
			d	&	a	&	c
		\end{array}
	\end{equation}
	\pause
	$a$ loses against $b$; $b$ against $d$; $c$ against $b$; $d$ against $c$
	\begin{itemize}
		\item Dodgson’s method (1876): candidates “closest” to being Condorcet winners {\tiny (in nb of swaps)}
	\end{itemize}
\end{frame}

\section{Analyzing voting rules}
\subsection[Problem]{Problems do arise}
\begin{frame}
	\frametitle{\subsecname}
	\begin{itemize}
		\item Example of a problem: Borda does not satisfy Condorcet’s principle
		\item Even when knowing the precise rule, we may overlook some problems of it
		\item We would like some way of finding out systematically whether a rule is “well-behaved”
	\end{itemize}
\end{frame}

\subsection{Axiomatics}
\begin{frame}
	\frametitle{\subsecname}
	\begin{quote}
		Rather than dream up a multitude of arbitration schemes and determine whether or not each withstands the best of plausibility in a host of special cases, let us invert the procedure. Let us examine our subjective intuition of fairness and formulate this as a set of precise desiderata that any acceptable arbitration scheme must fulfil. Once these desiderata are formalized as axioms, then the problem is reduced to a mathematical investigation of the existence of and characterization of arbitration schemes which satisfy the axioms.
	\end{quote}
	\citet[p. 121]{luce_games_1957}\par
	Some example of axioms? \pause Condorcet’s principle!
\end{frame}

\begin{frame}
	\frametitle{What’s an axiom?}
	\begin{itemize}
		\item An axiom (for us) is a principle
		\item Expressed formally
		\item That dictates some behavior of a voting rule
		\item In some conditions
		\item Usually seen as something to be satisfied
		\item Ideally, some union of some such axioms define exactly one rule
		\item Some axioms can be shown to be incompatible
	\end{itemize}
	\begin{example}[Condorcet’s principle]
		We ought to take the Condorcet winner as sole winner if it exists.
	\end{example}
\end{frame}

\subsection{Unanimity}
\begin{frame}
	\frametitle{\subsecname}
	\begin{definition}[Unanimity]
		We ought to select as winner someone who has no unanimously preferred alternative
	\end{definition}
	\begin{equation}
		\prof =
		\begin{array}{rrr}
			a	&	a	&	b\\
			b	&	b	&	c\\
			c	&	c	&	a
		\end{array}
	\end{equation}
	Constraint? \pause Do not take $c$ as $b$ is unanimously preferred to it.	\pause
	\begin{equation}
		\prof =
		\begin{array}{rrr}
			a	&	a	&	b\\
			b	&	c	&	c\\
			c	&	b	&	a
		\end{array}
	\end{equation}
	Constraint? \pause No constraint.
\end{frame}

\subsection[IIF]{Independence of Infeasible Alternatives}
\begin{frame}[fragile]
	\frametitle{Social choice correspondence}
%	\begin{description}[Social choice correspondence]
%		\item[Alternatives] $\allalts = \set{a, b, c, d, \ldots}$
%		\item[Feasible alternatives] $\emptyset \subset F \subseteq \allalts$
%		\item[Voters] $\allvoters = \set{1, 2, \ldots}$
%		\item[Profile] function $\prof$ from $\allvoters$ to orders on $\allalts$
%		\item[Social choice correspondence] function $f$ mapping each $\prof$ \emph{and} $\feasalts$ to winners $\emptyset \subset \alts \subseteq \feasalts$.
%	\end{description}
	\begin{itemize}
		\item Some alternatives might be infeasible
		\item Preferences over $\allalts$, choose among $\feasalts$
		\item $f$ a function that, given a profile and a set of feasible alternatives, chooses a subset of winners
		\item $f(\prof, \feasalts) \subseteq \feasalts$
	\end{itemize}
	\vfill
	\begin{center}
		\begin{tikzpicture}
			\path node (feas) {$\feasalts = \{a, c\}$};
			\path (feas.south) node[anchor=north, profile matrix] (profile) {
				\prof(1)&
				\prof(2)
				\\
				| (profile11) | a&
				b
				\\
				b&
				a
				\\
				c&
				| (profile32) | c
				\\
			};
			\path ($(profile.south west)!.5!(profile.south east)$) ++ (0, -5mm) node {$\prof$};
			
			\path node[draw, rectangle, fit=(profile11) (profile32), outer xsep=2mm, outer ysep=1mm] (justprofile) {};
			\path (justprofile.east) ++ (2.5cm, 0) node[inner sep=0] (winners) {\mbox{} $\alts = \set{a}$};
			\path[draw, ->] (justprofile.east) to[bend left=45] node[anchor=south] {$f$} (winners.west);
			\path[draw, ->] (feas.east) to[bend left=20] (winners.west);

			\path[draw, decorate, decoration={brace, mirror}] (justprofile.south west) -- (justprofile.south east);
		\end{tikzpicture}
	\end{center}
\end{frame}

\begin{frame}[fragile]
	\frametitle{\subsecname}
	\begin{definition}[\subsecname]
		When preferences change only about unfeasible alternatives, we ought not to change our decision
	\end{definition}
	\vspace{-1.8pt}
	\begin{example}
		\begin{equation}
			\begin{array}{l@{\hskip 0in}ll}
				\prof =
				&\begin{array}{rrr}
					b	&	c	&	d\\
					c	&	d	&	b\\
					a	&	b	&	a\\
					d	&	a	&	c
				\end{array},
				& f(\prof, \{b, c, d\}) = \set{b},\\[3em]
				\prof' =
				&\begin{array}{rrr}
					b	&	a	&	d\\
					c	&	c	&	a\\
					d	&	d	&	b\\
					a	&	b	&	c
				\end{array},
				& f(\prof', \{b, c, d\})\text{ ?}
				\pause
				\set{b}
			\end{array}
		\end{equation}
	\end{example}
\end{frame}

\subsection[ACA]{Arrow’s choice axiom}
\begin{frame}
	\frametitle{\subsecname}
	\begin{definition}[\subsecname]
		When feasible alternatives are restricted, with no change of preference, we ought not to change our decision if possible
	\end{definition}
	\vspace{-1.8pt}
	\begin{example}
		\begin{equation}
			\prof =
			\begin{array}{rrr}
				b	&	c	&	d\\
				c	&	d	&	b\\
				a	&	b	&	a\\
				d	&	a	&	c
			\end{array},
			f(\prof, \{b, c, d\}) = \set{b, d},
			f(\prof, \{b, c\})\text{ ?}
			\pause
			\set{b}
		\end{equation}
	\end{example}
\end{frame}

\subsection[Arrow]{Arrow’s impossibility theorem}
\begin{frame}
	\frametitle{Dictatorial social choice correspondences}
	\begin{definition}[Dictatorial social choice correspondence]
		$f$ is \emph{dictatorial} for a given voter iff $f$ always selects her preferred alternative among the feasible ones
	\end{definition}
	\vspace{-1.8pt}
	\begin{example}
		\begin{equation}
			\begin{array}{l@{\hskip 0in}ll}
				\prof =
				&\begin{array}{rrr}
					b	&	c	&	d\\
					c	&	d	&	b\\
					a	&	b	&	a\\
					d	&	a	&	c
				\end{array},
				& f(\prof, \{b, c, d\}) = \set{c},\\[3em]
				\prof' =
				&\begin{array}{rrr}
					b	&	a	&	d\\
					c	&	b	&	a\\
					d	&	d	&	b\\
					a	&	c	&	c
				\end{array},
				& f(\prof', \{b, c, d\})\text{ ?}
				\pause
				\set{b}
			\end{array}
		\end{equation}
	\end{example}
\end{frame}

\begin{frame}
	\frametitle{\subsecname}
	\begin{theorem}[\subsecname]
		When there are at least three alternatives, there is no social choice correspondence that satisfies Unanimity, Independence of Infeasible alternatives, Arrow’s Choice Axiom, and that is not dictatorial.
	\end{theorem}
\end{frame}

\section{Analyzing the Borda rule}
\subsection{Elementary profile}
\begin{frame}[fragile]
	\frametitle{\subsecname}
	
	Fix an arbitrary linear order $k$ on $\allalts$. Given $A \subseteq \allalts$, define $\pelem{A}$.
	\begin{minipage}{5.7cm}
		\begin{definition}[Elementary profile]
			\begin{tikzpicture}
				\path node[profile matrix, row 2/.style={nodes={draw, rectangle, minimum width=1.5em, minimum height=9ex}}, row 3/.style={nodes={draw, rectangle, minimum width=1.5em, minimum height=6ex}}] (profile) {
					\pelem{A}(1)	&\pelem{A}(2)\\
					| (profile-1-1) | 	& | (profile-1-2) | 	\\
					| (profile-2-1) | 	& | (profile-2-2) | 	\\
				};
				\path (profile-1-1.west) ++(-5mm, 0) node[anchor=east] (profile-1-1 expl) {$\restr{k}{A}$};
				\path (profile-1-2.east) ++(5mm, 0) node[anchor=west] (profile-1-2 expl) {$\restr{k^{-1}}{A}$};
				\path (profile-2-1.west) ++(-5mm, 0) node[anchor=east] (profile-2-1 expl) {$\restr{k}{\allalts \setminus A}$};
				\path (profile-2-2.east) ++(5mm, 0) node[anchor=west] (profile-2-2 expl) {$\restr{k^{-1}}{\allalts \setminus A}$};
				\path[draw, ->] (profile-1-1 expl.east) -- (profile-1-1.west);
				\path[draw, <-] (profile-1-2.east) -- (profile-1-2 expl.west);
				\path[draw, ->] (profile-2-1 expl.east) -- (profile-2-1.west);
				\path[draw, <-] (profile-2-2.east) -- (profile-2-2 expl.west);
				
				\path[draw, decorate, decoration={brace, mirror}] (profile.south west) -- (profile.south east);
				\path (profile.south) ++(0, -5mm) node {$\pelem{A}$};
			\end{tikzpicture}
		\end{definition}
	\end{minipage}\hspace{2em}%
	\begin{minipage}{\columnwidth - 5.7cm - 2em}
		\begin{example}
			\begin{tikzpicture}
				\path node[profile matrix] (profile) {
					\pelem{\Set{a, b, c}}(1)	&\pelem{\Set{a, b, c}}(2)\\
					| (profile-1-1) | a	& | (profile-1-2) | c	\\
					b	&b	\\
					| (profile-3-1) | c	& | (profile-3-2) | a	\\[2ex]
					| (profile-4-1) | d	& | (profile-4-2) | e	\\
					| (profile-5-1) | e	& | (profile-5-2) | d	\\
				};
				\path node[draw, rectangle, fit=(profile-1-1) (profile-3-1)] {};
				\path node[draw, rectangle, fit=(profile-4-1) (profile-5-1)] {};
				\path node[draw, rectangle, fit=(profile-1-2) (profile-3-2)] {};
				\path node[draw, rectangle, fit=(profile-4-2) (profile-5-2)] {};
				
				\path (profile.south west) ++ (1mm, -1mm) coordinate (profile-bottom-left);
				\path (profile.south east) ++ (-1mm, -1mm) coordinate (profile-bottom-right);
				\path[draw, decorate, decoration={brace, mirror}] (profile-bottom-left) -- (profile-bottom-right);
				\path (profile.south) ++(0, -5mm) node {$\pelem{\Set{a, b, c}}$};
			\end{tikzpicture}
		\end{example}
	\end{minipage}
\end{frame}

\begin{frame}
	\frametitle{Elementary profile: axiom}
	\begin{definition}[Axiom about elementary profiles]
		$f(\pelem{A}) = A$
	\end{definition}
\end{frame}

\subsection{Cyclic profile}
\begin{frame}
	\frametitle{\subsecname}
	\begin{definition}[Axiom about cyclic profiles]
		When given a cyclic profile, we ought to select all winners as ex-æquo
	\end{definition}
	\begin{example}
		\begin{equation}
			f\left(
			\begin{array}{cccc}
				a&b&c&d\\
				b&c&d&a\\
				c&d&a&b\\
				d&a&b&c\\
			\end{array}\right) = \allalts
		\end{equation}
	\end{example}
\end{frame}

\subsection{Cancellation}
\begin{frame}
	\frametitle{\subsecname}
	\begin{definition}[Cancellation]
		When all pairs of alternatives $(a, b)$ in a profile are such that $a$ is preferred to $b$ as many times as $b$ to $a$, we ought to select all winners as ex-æquo
	\end{definition}
	\begin{example}
		\begin{equation}
			f\left(%
			\begin{array}{rrrr}
				a&b&c&c\\
				b&a&a&b\\
				c&c&b&a\\
			\end{array}\right) = \allalts
		\end{equation}
	\end{example}
\end{frame}

\subsection{Reinforcement}
\begin{frame}
	\frametitle{\subsecname}
	\begin{definition}[\subsecname]
		When joining two sets of voters, exactly those winners that each set accepts should be selected, if possible
	\end{definition}
	\begin{example}
		\begin{equation}
			\prof_1 =
			\begin{array}{cc}
				a&b\\
				b&a\\
				c&c\\
			\end{array},
			A_1 = \set{a, b},
			\prof_2 =
			\begin{array}{ccc}
				a&b&a\\
				b&a&c\\
				c&c&b\\
			\end{array},
			A_2 = \set{a},
		\end{equation}
		\begin{equation}
			\prof =
			\begin{array}{ccccc}
				a&b&a&b&a\\
				b&a&b&a&c\\
				c&c&c&c&b\\
			\end{array}.
			\text{ Winners? }
			\pause
			\set{a}
		\end{equation}
	\end{example}
\end{frame}

\subsection{Characterisation of Borda}
\begin{frame}
	\frametitle{\subsecname}
	\begin{theorem}
		The Borda rule is the only voting rule that satisfies Elementary and Cyclic profile axioms, Cancellation and Reinforcement.
	\end{theorem}
	(We proved this variant of a characterisation given by \citet{young_axiomatization_1974})
\end{frame}

\subsection[Against Condorcet]{Fishburn-against-Condorcet argument}
\begin{frame}
	\frametitle{\subsecname}
	\begin{minipage}{5.4cm}
		\citet[p. 544]{fishburn_paradoxes_1974} argument against the Condorcet principle (see also \url{http://rangevoting.org/FishburnAntiC.html}).
		\begin{block}{Condorcet winner}
			$w \text{ VS } \mu, \mu \in \set{a, \ldots, h}\text{? }\uncover<2->{51/101}$
		\end{block}
	\end{minipage}\hfill%
	\mbox{%
		\small%
		$%
			\begin{array}{lrrrrrr}
				&\multicolumn{6}{c}{\text{nb voters}}\\
			\cmidrule{2-7}
					&31	&19	&10	&10	&10	&21	\\
			\midrule
				1	&a	&a	&f	&g	&h	&h	\\
				2	&b	&b	&w	&w	&w	&g	\\
				3	&c	&c	&a	&a	&a	&f	\\
				4	&d	&d	&h	&h	&f	&w	\\
				5	&e	&e	&g	&f	&g	&a	\\
				6	&w	&f	&e	&e	&e	&e	\\
				7	&g	&g	&d	&d	&d	&d	\\
				8	&h	&h	&c	&c	&c	&c	\\
				9	&f	&w	&b	&b	&b	&b	\\
			\end{array}
		$%
	}
	\centering
	$
		\begin{array}{lrrrrrrrrr}
			&\multicolumn{9}{c}{\text{ranks}}\\
		\cmidrule{2-10}
			&1	&≤2	&≤3	&≤4	&≤5	&≤6	&≤7	&≤8	&≤9	\\
		\midrule
		w	&0	& 30	& 30	& 51	& 51	& 82	& 82	& 82	&101	\\
		a	&50	& 50	& 80	& 80	& 101	& 101	& 101	&101	&101	\\
		\end{array}
	$
\end{frame}

\section{Our contribution}
\subsection{Goal}
\begin{frame}
	\frametitle{\subsecname}
	\begin{itemize}
		\item Better understanding of rules
		\item Illustrate reasoning of ≠ rules
		\item On examples
		\item On the basis of axioms
		\item Long-term goal: compare rules by letting them argue
	\end{itemize}
\end{frame}

\subsection{Approach}
\begin{frame}
	\frametitle{\subsecname}
	\begin{itemize}
		\item Express some axioms in a simple formal language
		\item Compute automatically their applications on examples
		\item Apply to the Borda rule
	\end{itemize}
\end{frame}

\subsection{Formal language}
\begin{frame}
	\frametitle{\subsecname}
	\begin{itemize}
		\item We use propositional logic
		\item A restricted language for easier deductions
		\item Contains $¬, ∨,∧, →$ but not $\forall, \exists$
		\item We translate axioms into this language
	\end{itemize}
\end{frame}

\subsection[Aim]{What we aim at}
\begin{frame}
	\frametitle{Example profile}
	\begin{block}{Who should win?}
		\begin{equation}
			\prof =
			\begin{array}{rrr}
				a&a&c\\
				b&b&b\\
				c&c&a\\
			\end{array}
		\end{equation}
	\end{block}
	\begin{itemize}
		\item Veto rule chooses $b$
		\item Borda rule chooses $a$
	\end{itemize}
\end{frame}

\begin{frame}
	\frametitle{\subsecname}
	\vspace{-13pt}
	\begin{equation}
		\prof =
		\begin{array}{>{\color{red}}r>{\color{darkgreen}}r>{\color{darkgreen}}r}
			a&a&c\\
			b&b&b\\
			c&c&a\\
		\end{array}
	\end{equation}
	\begin{tabularx}{\linewidth}{@{}>{\bfseries}lX>{\color{gray}[}r<{]}@{}}
		System: & Take the \textcolor{red}{\textit{red subprofile}}. Here, \textcolor{red}{\textit{$a$ should win}}, right? & unanimity \\
		User: & Obviously! \\
		System: & Now consider the \textcolor{darkgreen}{\textit{green subprofile}}. For symmetry reasons, there should be a \emph{three-way tie}, right? & cancellation \\
		User: & Sounds reasonable. \\
		System: &  So, as there was a three-way tie for the green part, the red part should decide the overall winner, right? & reinforcement \\
		User: & Yes. \\
		System: & \multicolumn{2}{l}{To summarise, you agree that $a$ should win.}
	\end{tabularx}
\end{frame}

\subsection{Provided}
\begin{frame}
	\frametitle{\subsecname}
	\begin{itemize}
		\item Our system can automatically explain the outcome of Borda
		\item Starting from any profile
		\item Requires solving a system of equation to find right intermediate profiles
		\item Does not use natural language (for the moment)
	\end{itemize}
\end{frame}

\subsection{Conclusion}
\begin{frame}
	\frametitle{\subsecname}
	\begin{itemize}
		\item No perfect voting rule
		\item Strengths can be expressed using axioms
		\item We use a specific language to express some of these axioms
		\item We can then instanciate concrete arguments (example-based)
		\item May render some arguments in the specialized literature accessible to non experts
		\item Extensions may permit to \emph{debate} about voting rules
		\item Provides a way to study appreciation of arguments
	\end{itemize}
\end{frame}

\begin{frame}[plain]
	\addtocounter{framenumber}{-1}
	\begin{center}
		\huge
		\textit{Thank you for your attention!}
	\end{center}
\end{frame}

\appendix
\AtBeginSection{
}

\begin{frame}[allowframebreaks]
	\frametitle{Bibliography}
	\def\newblock{\hskip .11em plus .33em minus .07em}
 	\bibliography{zotero}
\end{frame}

\section{Analyzing rules}
\subsection{Naïve attempt}
\begin{frame}
	\frametitle{\subsecname}
	⇒ Examine outcome of a rule on every possible profile?
	\begin{enumerate}
		\item We may not see what the problem is on a given example
		\item Too time consuming
	\end{enumerate}
	4 alternatives, 6 voters: at least $(4!)^6/4!/6!$ profiles $\eqsim \si{10000}$ {\tiny deducing obvious symetries for anonymity and neutrality}
% 4 voters, 4 alts: nb profiles = (4!)^4 but divided by approx 4! for permutations of alts, and divided by (very approx) 4! for permutations of voters, thus only 24×24 = 576.
\end{frame}

\section{Axioms}
\subsection{Anonymity}
\begin{frame}
	\frametitle{\subsecname}
	\begin{definition}[Anonymity]
		The identity of the voters should not matter
	\end{definition}
	\begin{equation}
		\prof_1 =
		\begin{array}{rrr}
			1&2&3\\
			a&b&b\\
			b&a&c\\
			c&c&a\\
		\end{array},
		\prof_2 =
		\begin{array}{ccc}
			3&1&2\\
			a&b&b\\
			b&a&c\\
			c&c&a\\
		\end{array}
	\end{equation}
	Constraint? \pause Same winners.
\end{frame}

\subsection{Cyclic profiles}
\begin{frame}
	\frametitle{\subsecname}
	Given $S$ a complete cycle in $\allalts$, define $\pcycl{S}$.
	\begin{definition}[Cyclic profile]
		$\pcycl{S}$ is the profile composed by all $\card{\allalts}$ possible linearizations of $S$ as preference orderings. 
	\end{definition}
	\begin{example}
		\begin{equation}
			\pcycl{\left<a, b, c, d\right>} =
			\begin{array}{cccc}
				a&b&c&d\\
				b&c&d&a\\
				c&d&a&b\\
				d&a&b&c\\
			\end{array}.
		\end{equation}
	\end{example}
\end{frame}

\subsection{Reinforcement}
\begin{frame}
	\frametitle{Reinforcement l-axiom}
	Consider $\prof_1$, $\prof_2$,
	\begin{itemize}
		\item having winners $A_1$, $A_2$,
		\item with $A_1 \cap A_2 ≠ \emptyset$;
	\end{itemize}
	then winners in $\prof_1 + \prof_2$ must be $A_1 \cap A_2$.
	\begin{definition}[\laxiom{Reinf}]
		For each $\prof_1, \prof_2, A_1, A_2 \subseteq \allalts, A_1 \cap A_2 ≠ \emptyset$:
		\begin{equation}
			(\lato[\prof_1][A_1] ∧ \lato[\prof_2][A_2]) → \lato[\prof_1 + \prof_2][A_1 \cap A_2].
		\end{equation}
	\end{definition}
\end{frame}

\section{Example using Borda}
\begin{frame}
	\frametitle{An example}
	
	Consider $\allalts=\set{a, b, c, d}$ and a profile $\prof$ defined as:
	\begin{equation}
		\prof =
		\begin{array}{cc}
			a	&	c\\
			b	&	b\\
			d	&	a\\
			c	&	d
		\end{array}.
	\end{equation}
	\pause
	We want to justify that $f_B(\prof)=\Set{a, b}$.
\end{frame}

\begin{frame}
	\frametitle{Sketch}
	
	\begin{itemize}
		\item Consider any $\prof' = q_1 \pelem{a, b} + q_2 \pelem{a, b, c} + \sum_{S \in \mathcal{S}} q_S \pcycl{S}$, $q_1, q_2, q_S \in \N, \mathcal{S}$ some set of cycles.
		\item In $\prof'$, $W=\set{a, b}$ must win.
		\item Find $k \in \N$ such that $\overline{k \prof} + \prof'$ cancel.
		\item Then $k \prof$ has winners $W$. (Skipping details.)
		\item Then $\prof$ has winners $W$.
	\end{itemize}
	Our task: find $\prof'$ a combination of elementary and cyclic profiles such that $\overline{k \prof} + \prof'$ cancel.
	
	Good news: this is always possible.
\end{frame}

\begin{frame}
	\frametitle{Application on the example}
	
	Define $\prof' = \pelem{a, b} + 2 \pelem{a, b, c} + \pcycl{\left<c, b, a, d\right>} + \pcycl{\left<b, d, c, a\right>}$.
	\begin{enumerate}
		\item $\lato[\pelem{a, b}][\set{a, b}]$ (\laxiom{elem})
		\item $\lato[\pelem{a, b, c}][\set{a, b, c}]$ (\laxiom{elem})
		\item $\lato[\pcycl{\left<c, b, a, d\right>}][\allalts]$ (\laxiom{cycl})
		\item $\lato[\pcycl{\left<b, d, c, a\right>}][\allalts]$ (\laxiom{cycl})
		\item $\lato[\prof'][\set{a, b}]$ (\laxiom{reinf}, 1, 2, 3, 4)
		\item $\lato[4 \prof + \overline{4 \prof}][\allalts]$ (\laxiom{canc})
		\item $\lato[4 \prof + \overline{4 \prof} + \prof'][\set{a, b}]$ (\laxiom{reinf}, 5, 6)
		\item $\lato[\overline{4 \prof} + \prof'][\allalts]$ (\laxiom{canc})
		\item $\lato[4 \prof][\set{a, b}]$ (\laxiom{reinf}, 7, 8)
		\item $\lato[\prof][\set{a, b}]$ (\laxiom{reinf}, 9)
	\end{enumerate}
\end{frame}

\section{License}
\begin{frame}
	\frametitle{License}
	This presentation, and the associated \LaTeX{} code, are published under the \href{http://opensource.org/licenses/MIT}{MIT license}. Feel free to reuse (parts of) the presentation, under condition that you cite the author.
	
	Credits are to be given to \href{http://www.lamsade.dauphine.fr/~ocailloux/}{Olivier Cailloux}, Université Paris-Dauphine.
\end{frame}
\end{document}

\section{Frame template}
\subsection{}
\begin{frame}
	\frametitle{\subsecname}
	\begin{itemize}
		\item 
	\end{itemize}
\end{frame}

